\chapter{عملگرها و محاسبه}
\label{ch:operators}

\begin{goals}
\begin{itemize}
  \item عملگرهای حسابی و ترتیب انجام آن‌ها
  \item عملگرهای مقایسه و نتیجه‌ی منطقی آن‌ها
  \item ترکیب شرط‌ها با \fcode{و}، \fcode{یا} و \code{!}
  \item شکل‌های کوتاه مقداردهی
  \item چند محاسبه‌ی روزمره با همین ابزارها
\end{itemize}
\end{goals}

\section{عملگر چیست؟}

علامت‌هایی مثل \code{+} و \code{*} که با آن‌ها محاسبه می‌کنیم،
\textbf{عملگر} نام دارند و چیزهایی که روی آن‌ها کار می‌کنند
\textbf{عملوند}. در \code{2 + 3}، علامت \code{+} عملگر است و \code{2} و
\code{3} عملوندها. تا این‌جا با چند عملگر آشنا شده‌اید؛ در این فصل همه‌ی
عملگرهای پایه‌ی سلام را یک‌جا می‌بینیم.

\section{عملگرهای حسابی}

\begin{center}
\begin{tabular}{clc}
\toprule
عملگر & معنی & نمونه \\
\midrule
\code{+} & جمع & \code{7 + 2} برابر \code{9} \\
\code{-} & تفریق & \code{7 - 2} برابر \code{5} \\
\code{*} & ضرب & \code{7 * 2} برابر \code{14} \\
\code{/} & تقسیم & \code{7 / 2} برابر \code{3} \\
\code{\%} & باقی‌مانده & \code{7 \% 2} برابر \code{1} \\
\code{\textasciicircum\textasciicircum} & توان & \code{2 \textasciicircum\textasciicircum 3} برابر \code{8} \\
\bottomrule
\end{tabular}
\end{center}

\program{عملگرهای حسابی}
\begin{salam}
تابع اصلی:
    چاپ 7 + 2, 7 - 2, 7 * 2
    چاپ 7 / 2, 7 % 2
    چاپ 2 ^^ 3, 2 ^^ 10
    چاپ 10 - 15
تمام
\end{salam}

\begin{outputfa}
9 5 14
3 1
8 1024
-5
\end{outputfa}

عملگر توان، \code{\textasciicircum\textasciicircum}، همیشه نتیجه‌ی
اعشاری می‌دهد، حتی اگر هر دو عملوند صحیح باشند. علامت منفی را می‌توان
پیش از عدد یا متغیر گذاشت تا قرینه‌اش را بدهد: \fcode{-سن}.

\section{ترتیب انجام عملگرها}

وقتی چند عملگر در یک عبارت باشند، سلام همان قاعده‌ای را رعایت می‌کند که در
ریاضی مدرسه یاد گرفته‌اید: اول توان، بعد ضرب و تقسیم و باقی‌مانده، بعد جمع
و تفریق. عملگرهای هم‌رتبه از راست به چپ نوشتار انگلیسی، یعنی به ترتیبی
که نوشته شده‌اند، انجام می‌شوند:

\program{ترتیب عملگرها}
\begin{salam}
تابع اصلی:
    چاپ 2 + 3 * 4
    چاپ (2 + 3) * 4
    چاپ 10 - 4 - 3
    چاپ 100 / 10 / 2
    چاپ 2 * 3 ^^ 2
تمام
\end{salam}

\begin{outputfa}
14
20
3
5
18
\end{outputfa}

\code{2 + 3 * 4} برابر \code{14} شد نه \code{20}، چون ضرب زودتر انجام
می‌شود. \code{10 - 4 - 3} برابر \code{3} است، چون اول \code{10 - 4} حساب
می‌شود و بعد \code{6 - 3}. هر جا شک دارید یا می‌خواهید ترتیب دیگری
بخواهید، پرانتز بگذارید. پرانتز اضافی هرگز ضرری ندارد و اغلب خواندن را
راحت‌تر می‌کند.

\begin{tip}
عبارت‌های خیلی طولانی را به چند متغیر با نام‌های گویا بشکنید. به جای
\fcode{چاپ (الف + ب) * ج / 100 - د}، بنویسید \fcode{جمع \lr{:=} الف + ب}
و سپس \fcode{تخفیف \lr{:=} جمع * ج / 100} و در آخر
\fcode{چاپ تخفیف - د}. خواننده (و خودتان یک هفته بعد) از شما ممنون
خواهد شد.
\end{tip}

\section{کارهایی که با تقسیم و باقی‌مانده می‌شود کرد}

تقسیم صحیح و باقی‌مانده، با هم، ابزاری قدرتمند برای «خرد کردن» عددها
هستند. چند نمونه:

\program{تبدیل ثانیه به ساعت و دقیقه}
\begin{salam}
تابع اصلی:
    مجموع ثانیه := 3725
    ساعت := مجموع ثانیه / 3600
    دقیقه := (مجموع ثانیه % 3600) / 60
    ثانیه := مجموع ثانیه % 60
    چاپ ساعت, "ساعت و", دقیقه, "دقیقه و", ثانیه, "ثانیه"
تمام
\end{salam}

\begin{outputfa}
1 ساعت و 2 دقیقه و 5 ثانیه
\end{outputfa}

هر ساعت ۳۶۰۰ ثانیه است، پس تقسیم صحیح بر ۳۶۰۰ تعداد ساعت‌های کامل را
می‌دهد. باقی‌مانده‌ی همین تقسیم، ثانیه‌هایی است که یک ساعت کامل نساخته‌اند و
همان را بر ۶۰ تقسیم می‌کنیم تا دقیقه‌ها به دست بیاید.

\program{رقم‌های یک عدد}
\begin{salam}
تابع اصلی:
    عدد := 472
    چاپ "یکان:", عدد % 10
    چاپ "دهگان:", عدد / 10 % 10
    چاپ "صدگان:", عدد / 100
تمام
\end{salam}

\begin{outputfa}
یکان: 2
دهگان: 7
صدگان: 4
\end{outputfa}

باقی‌مانده‌ی تقسیم بر ۱۰، همیشه رقم یکان است. تقسیم بر ۱۰ یک رقم را از
سمت راست می‌اندازد. با ترکیب این دو می‌توان به هر رقمی رسید.

\program{درصد و تخفیف}
\begin{salam}
تابع اصلی:
    قیمت := 250000
    درصد تخفیف := 15
    مبلغ تخفیف := قیمت * درصد تخفیف / 100
    چاپ "تخفیف:", مبلغ تخفیف
    چاپ "قیمت نهایی:", قیمت - مبلغ تخفیف
تمام
\end{salam}

\begin{outputfa}
تخفیف: 37500
قیمت نهایی: 212500
\end{outputfa}

دقت کنید که اول ضرب کردیم و بعد تقسیم. اگر می‌نوشتیم
\fcode{قیمت * (درصد تخفیف / 100)}، تقسیم صحیح ۱۵ بر ۱۰۰ صفر می‌شد و کل
تخفیف از بین می‌رفت. با عددهای صحیح، تا جای ممکن تقسیم را آخر انجام دهید.

\section{شکل‌های کوتاه مقداردهی}

در فصل پیش با \code{+=} و \code{++} آشنا شدید. همه‌ی عملگرهای حسابی شکل
کوتاه دارند:

\begin{center}
\begin{tabular}{ll}
\toprule
شکل کوتاه & معادل \\
\midrule
\fcode{الف \lr{+=} 5} & \fcode{الف = الف \lr{+} 5} \\
\fcode{الف \lr{-=} 5} & \fcode{الف = الف \lr{-} 5} \\
\fcode{الف \lr{*=} 5} & \fcode{الف = الف \lr{*} 5} \\
\fcode{الف \lr{/=} 5} & \fcode{الف = الف \lr{/} 5} \\
\fcode{الف \lr{\%=} 5} & \fcode{الف = الف \lr{\%} 5} \\
\fcode{الف\lr{++}} & \fcode{الف = الف \lr{+} 1} \\
\fcode{الف\lr{--}} & \fcode{الف = الف \lr{-} 1} \\
\bottomrule
\end{tabular}
\end{center}

\program{شکل‌های کوتاه}
\begin{salam}
تابع اصلی:
    متغیر امتیاز := 100
    امتیاز += 20
    چاپ امتیاز
    امتیاز -= 50
    چاپ امتیاز
    امتیاز *= 2
    چاپ امتیاز
    امتیاز /= 4
    چاپ امتیاز
    امتیاز++
    چاپ امتیاز
    امتیاز--
    چاپ امتیاز
تمام
\end{salam}

\begin{outputfa}
120
70
140
35
36
35
\end{outputfa}

\section{عملگرهای مقایسه}

عملگرهای مقایسه دو مقدار را با هم می‌سنجند و نتیجه‌ی منطقی می‌دهند:
\fcode{درست} یا \fcode{نادرست}.

\begin{center}
\begin{tabular}{cl}
\toprule
عملگر & معنی \\
\midrule
\code{==} & برابر است با \\
\code{!=} & برابر نیست با \\
\code{<} & کوچک‌تر از \\
\code{>} & بزرگ‌تر از \\
\code{<=} & کوچک‌تر یا مساوی \\
\code{>=} & بزرگ‌تر یا مساوی \\
\bottomrule
\end{tabular}
\end{center}

\program{مقایسه}
\begin{salam}
تابع اصلی:
    سن := 20
    چاپ سن == 20
    چاپ سن != 20
    چاپ سن < 18
    چاپ سن >= 18
    چاپ "سیب" == "سیب"
    چاپ "سیب" == "گلابی"
تمام
\end{salam}

\begin{outputfa}
true
false
false
true
true
false
\end{outputfa}

\begin{warn}
اشتباه \code{=} با \code{==} یکی از رایج‌ترین اشتباه‌های همه‌ی
برنامه‌نویسان جهان است. \code{=} مقدار می‌دهد، \code{==} می‌پرسد «آیا
برابرند؟». اگر جایی که باید مقایسه کنید \code{=} بنویسید، سلام خطا
می‌دهد، که خبر خوبی است.
\end{warn}

نتیجه‌ی مقایسه یک مقدار منطقی معمولی است و می‌توان آن را در متغیری نگه
داشت. این کار وقتی شرطی پیچیده است یا چند بار استفاده می‌شود، برنامه را
خواناتر می‌کند:

\begin{salam}
تابع اصلی:
    دما := 38
    هوا گرم است := دما > 30
    چاپ "هوا گرم است؟", هوا گرم است
تمام
\end{salam}

\begin{outputfa}
هوا گرم است؟ true
\end{outputfa}

\section{ترکیب شرط‌ها: و، یا، نقیض}

بسیاری از تصمیم‌ها به بیش از یک شرط بستگی دارند: «اگر بزرگسال است \emph{و}
بلیت دارد»، «اگر جمعه است \emph{یا} تعطیل رسمی است». سلام برای ترکیب
شرط‌ها سه عملگر دارد که دو تای آن‌ها واژه‌های فارسی‌اند:

\begin{center}
\begin{tabular}{cll}
\toprule
عملگر & معنی & نتیجه درست است وقتی که \\
\midrule
\fcode{و} & و & هر دو طرف درست باشند \\
\fcode{یا} & یا & دست‌کم یکی از دو طرف درست باشد \\
\code{!} & نقیض (نه) & طرف راستش نادرست باشد \\
\bottomrule
\end{tabular}
\end{center}

\program{ترکیب شرط‌ها}
\begin{salam}
تابع اصلی:
    سن := 20
    بلیت دارد := نادرست
    چاپ سن >= 18 و بلیت دارد
    چاپ سن >= 18 یا بلیت دارد
    چاپ !بلیت دارد
    چاپ سن > 12 و سن < 65
تمام
\end{salam}

\begin{outputfa}
false
true
true
true
\end{outputfa}

خط آخر را ببینید: برای پرسیدن «آیا سن میان ۱۲ و ۶۵ است؟» باید دو مقایسه‌ی
جداگانه را با \fcode{و} به هم وصل کرد. نوشتن \fcode{12 < سن < 65}، آن‌طور
که در ریاضی می‌نویسیم، در سلام معنی ندارد.

\begin{note}
به جای \fcode{و} و \fcode{یا} می‌توانید \code{\&\&} و \code{||} هم
بنویسید؛ این‌ها علامت‌هایی هستند که در بسیاری از زبان‌های دیگر به کار
می‌روند و سلام هر دو شکل را می‌پذیرد. در این کتاب شکل فارسی را به کار
می‌بریم، چون خواناتر است.
\end{note}

\begin{deep}[جدول درستی]
اگر بخواهید همه‌ی حالت‌های \fcode{و} و \fcode{یا} را یک‌جا ببینید،
جدول زیر کمکتان می‌کند. آن را حفظ نکنید؛ فقط با معنی روزمره‌ی «و» و «یا»
مقایسه‌اش کنید و می‌بینید که کاملاً همان است.
\begin{center}
\begin{tabular}{cccc}
\toprule
الف & ب & \fcode{الف و ب} & \fcode{الف یا ب} \\
\midrule
درست & درست & درست & درست \\
درست & نادرست & نادرست & درست \\
نادرست & درست & نادرست & درست \\
نادرست & نادرست & نادرست & نادرست \\
\bottomrule
\end{tabular}
\end{center}
تنها جایی که «یا»ی برنامه‌نویسی با «یا»ی روزمره فرق دارد، وقتی است که هر
دو طرف درست باشند: در گفت‌وگوی روزمره «چای یا قهوه؟» یعنی یکی از آن دو،
اما \fcode{یا} در سلام وقتی هر دو درست باشند هم درست است.
\end{deep}

\section{یک برنامه‌ی کامل‌تر}

فرض کنید می‌خواهید بدانید برای یک مهمانی چند کارتن نوشیدنی بخرید. هر
کارتن ۶ بطری دارد و هر مهمان ۲ بطری می‌نوشد:

\program{خرید برای مهمانی}
\begin{salam}
تابع اصلی:
    تعداد مهمان := 17
    بطری هرنفر := 2
    بطری کارتن := 6

    بطری لازم := تعداد مهمان * بطری هرنفر
    کارتن کامل := بطری لازم / بطری کارتن
    بطری اضافه := بطری لازم % بطری کارتن
    // اگر بطری اضافه داریم، یک کارتن دیگر لازم است
    کارتن لازم := کارتن کامل + (بطری اضافه > 0 ؟ 1 : 0)

    چاپ "بطری لازم:", بطری لازم
    چاپ "کارتن لازم:", کارتن لازم
    چاپ "بطری باقی‌مانده:", کارتن لازم * بطری کارتن - بطری لازم
تمام
\end{salam}

\begin{outputfa}
بطری لازم: 34
کارتن لازم: 6
بطری باقی‌مانده: 2
\end{outputfa}

در این برنامه یک عملگر تازه هم هست: \fcode{شرط ؟ الف : ب}. این عملگر
می‌گوید «اگر شرط درست است مقدار الف، وگرنه مقدار ب». این‌جا اگر بطری
اضافه‌ای مانده باشد، یک کارتن بیشتر می‌خریم. (شاید بپرسید چرا ننوشتیم
\fcode{بطری هر نفر}؛ پاسخ در فصل \ref{ch:variables} است: واژه‌ی \fcode{هر}
یکی از واژه‌های خود زبان است و نمی‌تواند بخشی از یک نام باشد.) در فصل بعد، که به تصمیم‌گیری
می‌پردازد، این عملگر را بیشتر می‌بینیم.

\begin{summary}
\begin{itemize}
  \item عملگرهای حسابی: \code{+ - * / \%} و توان
        \code{\textasciicircum\textasciicircum}. تقسیم صحیح اعشار را
        دور می‌ریزد و \code{\%} باقی‌مانده را می‌دهد.
  \item ترتیب: اول توان، بعد ضرب و تقسیم، بعد جمع و تفریق. با پرانتز
        ترتیب را عوض کنید.
  \item شکل‌های کوتاه: \code{+=}، \code{-=}، \code{*=}، \code{/=}،
        \code{++}، \code{--}.
  \item مقایسه: \code{== != < > <= >=} که نتیجه‌ی منطقی می‌دهند.
  \item ترکیب شرط‌ها با \fcode{و}، \fcode{یا} و \code{!}.
  \item \fcode{شرط ؟ الف : ب} یکی از دو مقدار را بر پایه‌ی شرط انتخاب
        می‌کند.
\end{itemize}
\end{summary}

\section*{تمرین‌ها}

\exercise بدون اجرا کردن، نتیجه‌ی هر عبارت را بگویید و سپس بررسی کنید:
\code{5 + 2 * 3}، \code{(5 + 2) * 3}، \code{20 / 3 * 3}، \code{20 \% 3},
\code{2 \textasciicircum\textasciicircum 2 \textasciicircum\textasciicircum 3}.

\exercise عددی چهاررقمی را در متغیری بگذارید و مجموع رقم‌هایش را با تقسیم
و باقی‌مانده حساب و چاپ کنید.

\exercise قیمت یک کالا و مبلغی که مشتری پرداخته را در دو متغیر بگذارید و
باقی‌مانده را به کمترین تعداد اسکناس ۱۰٬۰۰۰ تومانی، ۵٬۰۰۰ تومانی و
۱٬۰۰۰ تومانی تقسیم کنید و تعداد هر یک را چاپ کنید.

\exercise سال تولد کسی را در متغیری بگذارید و با یک عبارت منطقی چاپ کنید
که آیا او امسال (۱۴۰۵) به سن رأی دادن (۱۸ سال) رسیده است یا نه.

\exercise سه متغیر منطقی به نام‌های \fcode{باران می‌بارد}،
\fcode{چتر دارم} و \fcode{ماشین دارم} بسازید و عبارتی بنویسید که بگوید
آیا «بدون خیس شدن» می‌توان بیرون رفت: یعنی یا باران نمی‌بارد، یا چتر یا
ماشین دارم. برنامه را با چند ترکیب مختلف از مقدارها امتحان کنید.
